\section{第41届复赛第5题：Ioffe磁阱与蒸发冷却}
\subsection{题干}
四根无限长直导线穿过边长为 $a$、$b$ 的长方形四个顶点，电流大小均为 $I$，方向如图5a。以中心 $O$ 为原点，研究原点附近 $P(x,y,0)$ 的磁场；进而证明磁矩原子的势能，求磁矩与局域磁场反向时的约束力，并分析半径从 $R_1$ 缩至 $R_1/5$、密度保持不变的蒸发冷却过程。
\begin{figure}[htbp]\centering\begin{tikzpicture}[>=Stealth,scale=.88]
  \draw[densely dotted,thick] (-2,-1.45) rectangle (2,1.45);
  \draw[->,very thick] (-2.65,0)--(2.75,0) node[right] {$x$};
  \draw[->,very thick] (0,-2.05)--(0,2.15) node[above] {$y$};
  \node[below left] at (0,0) {$O$};
  \draw[fill=white,thick] (-2,1.45) circle(.22); \node at (-2,1.45) {$\odot$}; \node[above] at (-2,1.7) {$I$};
  \draw[fill=white,thick] (2,1.45) circle(.22); \node at (2,1.45) {$\otimes$}; \node[above] at (2,1.7) {$I$};
  \draw[fill=white,thick] (-2,-1.45) circle(.22); \node at (-2,-1.45) {$\otimes$}; \node[below] at (-2,-1.7) {$I$};
  \draw[fill=white,thick] (2,-1.45) circle(.22); \node at (2,-1.45) {$\odot$}; \node[below] at (2,-1.7) {$I$};
  \fill (0.7,.62) circle(1.5pt); \node[above right] at (.7,.62) {$P(x,y)$};
  \draw[<->] (-2,-1.75)--(2,-1.75) node[midway,below] {$a$};
  \draw[<->] (2.35,-1.45)--(2.35,1.45) node[midway,right] {$b$};
  \node[below,align=center] at (0,-2.48) {原题图5a：四根无限长直导线构成的 Ioffe 势阱};
\end{tikzpicture}
\caption{第41届复赛第5题原题图5a的TikZ重绘。}\end{figure}
\subsection{参考解答整理}
无限长直导线在距导线 $r$ 处产生 $B=\mu_0I/(2\pi r)$。由四根导线的对称性，$O$ 点磁场为零。对 $x\ll a$、$y\ll b$，答案给出原点附近的领头阶近似
\[B_x=-\frac{16\mu_0abI}{\pi(a^2+b^2)^2}x,\qquad B_y=\frac{16\mu_0abI}{\pi(a^2+b^2)^2}y,\]
\[B=\frac{16\mu_0abI}{\pi(a^2+b^2)^2}\sqrt{x^2+y^2}.\]
原子磁矩在磁场中的势能为 $W_m=-\boldsymbol S\cdot\boldsymbol B$。本题通过实验使磁矩始终与磁场反向，故以 $O$ 为势能零点时
\[W_m=\frac{16\mu_0abIS}{\pi(a^2+b^2)^2}\sqrt{x^2+y^2},\]
且约束力指向原点：
\[F_x=-\frac{16\mu_0abIS}{\pi(a^2+b^2)^2}\frac{x}{\sqrt{x^2+y^2}},\qquad F_y=-\frac{16\mu_0abIS}{\pi(a^2+b^2)^2}\frac{y}{\sqrt{x^2+y^2}}.\]

半径为 $R_1$ 的均匀原子团有 $N_1=n\pi R_1^2$；其总势能和平均势能分别为
\[E_1=\frac{32\mu_0abISn}{3(a^2+b^2)^2}R_1^3,\qquad \bar E_1=\frac{32\mu_0abIS}{3\pi(a^2+b^2)^2}R_1.\]
冷却后 $R_2=R_1/5$，故 $N_2=N_1/25$、$E_2=E_1/125$、$\bar E_2=\bar E_1/5$。下图为解析辅助图，不是原题图。
\begin{figure}[htbp]\centering\begin{tikzpicture}[>=Stealth,scale=.9]
  \draw[thick] (-2.25,0) circle(1.45); \fill[blue!12] (-2.25,0) circle(1.45);
  \draw[<->] (-3.7,-1.75)--(-.8,-1.75) node[midway,below] {$2R_1$};
  \node[below] at (-2.25,-2.12) {初态：$N_1=n\pi R_1^2$};
  \draw[->,very thick] (-.35,0)--(.8,0) node[midway,above] {蒸发冷却};
  \draw[thick] (2.25,0) circle(.29); \fill[blue!12] (2.25,0) circle(.29);
  \draw[<->] (1.96,-.6)--(2.54,-.6) node[midway,below] {$2R_1/5$};
  \node[below,align=center] at (2.25,-1.0) {冷却后：$N_2=N_1/25$\\$E_2=E_1/125$，$\bar E_2=\bar E_1/5$};
  \node[below] at (0,-2.72) {解析辅助图：半径收缩时的原子数与势能比例};
\end{tikzpicture}
\caption{第41届复赛第5题解析辅助图：蒸发冷却后的半径收缩与量的比例。}\end{figure}
\subsection{转写依据}
题干依据 \texttt{CPHO41-Q05\_transcription\_v1.md}；公式与分值依据答案第11至13页和 \texttt{CPHO41-Q05\_formula\_check\_v1.csv}。图5a由原题图重绘；蒸发冷却图为明确标注的解析辅助图。
